\chapter{字符串}

\section{KMP}

\cardmeta{基础必会}{单模式匹配、重叠出现、周期}{\(O(n+m)\)}{模式串为空时按题面单独处理}

\begin{lstlisting}[style=cpp]
vector<int> prefix_function(const string &p) {
    // pi[i]：p[0..i] 的最长真前后缀长度
    vector<int> pi(p.size());
    for (int i = 1; i < (int)p.size(); ++i) {
        int j = pi[i - 1];
        // 失配时沿着已经求出的 border 链跳，不回退 i
        while (j && p[i] != p[j]) j = pi[j - 1];
        if (p[i] == p[j]) ++j;
        pi[i] = j;
    }
    return pi;
}

vector<int> kmp(const string &text, const string &pat) {
    if (pat.empty()) return {};
    auto pi = prefix_function(pat);
    vector<int> pos;
    for (int i = 0, j = 0; i < (int)text.size(); ++i) {
        // j 表示当前已经匹配的 pat 前缀长度
        while (j && text[i] != pat[j]) j = pi[j - 1];
        if (text[i] == pat[j]) ++j;
        if (j == (int)pat.size()) {
            pos.push_back(i - j + 1);
            j = pi[j - 1]; // 保留重叠前缀
        }
    }
    return pos;
}
\end{lstlisting}

\begin{card}
\textbf{样例}\\
输入：文本 `abababa`，模式 `aba`。\\
输出：匹配起点 `0 2 4`，数量为 `3`。\\
`aaaaa` 与 `aaa` 的起点是 `0 1 2`，用来检查匹配成功后不能把 `j` 直接清零。
\end{card}

\section{Trie 与 AC 自动机}

\cardmeta{现场常用}{多个模式串、文本扫描、后缀模式}{建树 \(O(S\Sigma)\)，扫描 \(O(|T|)\)}{模式总长度为 \(S\)，只支持小写字母}

下面这份代码的 `out[u]` 是“以节点 \(u\) 结尾的模式串数量”。`build()` 后把缺失转移补齐，因此只能在插入完成后调用一次；多测必须重新构造对象。

\begin{lstlisting}[style=cpp]
struct AC {
    // go 是 Trie 转移；fail[u] 指向 u 的最长真后缀节点
    // terminal[u] 只统计恰好在 u 结束的模式串（允许重复）
    // output[u] 是沿 fail 链可匹配的模式串总数
    vector<array<int, 26>> go;
    vector<int> fail, terminal, output;

    AC() { init(); }
    void init() {
        go.clear(); fail.clear();
        terminal.clear(); output.clear();
        go.push_back({}); go[0].fill(0);
        fail.push_back(0);
        terminal.push_back(0);
        output.push_back(0);
    }
    int insert(const string &s) {
        int u = 0;
        for (char ch : s) {
            int c = ch - 'a';
            if (!go[u][c]) {
                // 首次经过这条字符边时创建 Trie 节点
                go[u][c] = go.size();
                go.push_back({});
                go.back().fill(0);
                fail.push_back(0);
                terminal.push_back(0);
                output.push_back(0);
            }
            u = go[u][c];
        }
        ++terminal[u];
        return u;
    }
    void build() {
        queue<int> q;
        // 根的 fail 为 0；第一层节点直接入队
        for (int c = 0; c < 26; ++c) {
            if (go[0][c]) q.push(go[0][c]);
        }
        while (!q.empty()) {
            int u = q.front(); q.pop();
            // output 汇总当前节点以及 fail 链上的终止模式串
            output[u] = terminal[u] + output[fail[u]];
            for (int c = 0; c < 26; ++c) {
                int &v = go[u][c];
                if (v) {
                    // 已存在的 Trie 边：由父节点的 fail 转移计算 fail[v]
                    fail[v] = go[fail[u]][c];
                    q.push(v);
                } else {
                    // 补全自动机转移，查询时无需 while 回退
                    v = go[fail[u]][c];
                }
            }
        }
        output[0] = terminal[0];
    }
    int query_total(const string &text) const {
        int u = 0, ans = 0;
        for (char ch : text) {
            // 每个字符只走一次自动机边，累计当前位置结尾的所有模式
            u = go[u][ch - 'a'];
            ans += output[u];
        }
        return ans;
    }
};
\end{lstlisting}

\begin{card}
\textbf{完整样例}\\
模式串：`a`、`ab`、`bab`；文本：`abab`。\\
每个模式串出现次数分别为：
\[
a:2,\qquad ab:2,\qquad bab:1.
\]
如果需要分别统计，把 `insert()` 返回的结束节点保存下来，扫描后在 fail 树上汇总文本经过次数；不能只读取 Trie 节点本身。
\end{card}

\begin{warnbox}
`output[u]` 版本适合统计文本中所有模式串的总出现次数。若题目要求每个模式串分别输出，必须另写 fail 树统计版本，不能把两种语义混用。
\end{warnbox}

\subsection{AC 自动机：分别统计每个模式串}

仍然复用上一个 \texttt{AC} 对象。扫描文本时给当前状态计数，再沿 fail 树从深到浅汇总；因此
“模式串互为后缀”时，短模式也会收到长模式经过节点的贡献。返回的 \texttt{pattern\_node[i]}
是第 \(i\) 个模式串插入结束的节点，重复模式串可以保留多个相同节点。

\begin{lstlisting}[style=cpp]
vector<int> count_each_pattern(const AC &ac, const string &text,
                               const vector<int> &pattern_node) {
    vector<int> pass(ac.go.size(), 0), order;
    int u = 0;
    for (char ch : text) {
        // build() 已补齐转移，所以扫描阶段不需要沿 fail while 回退。
        u = ac.go[u][ch - 'a'];
        ++pass[u];
    }

    // fail 边反向后是一棵树；反向遍历保证先汇总深层节点。
    vector<vector<int>> fail_tree(ac.go.size());
    for (int v = 1; v < (int)ac.go.size(); ++v)
        fail_tree[ac.fail[v]].push_back(v);
    vector<int> st = {0};
    while (!st.empty()) {
        int x = st.back(); st.pop_back();
        order.push_back(x);
        for (int v : fail_tree[x]) st.push_back(v);
    }
    reverse(order.begin(), order.end());
    for (int x : order)
        if (x) pass[ac.fail[x]] += pass[x];

    vector<int> answer;
    answer.reserve(pattern_node.size());
    for (int node : pattern_node) answer.push_back(pass[node]);
    return answer;
}
\end{lstlisting}

\begin{card}
\textbf{完整输入输出：}
\begin{lstlisting}
模式串：a, ab, bab
文本：abab
每个模式串出现次数：2 2 1
\end{lstlisting}
\textbf{四组边界：}模式集合为空；重复模式 `a,a`；后缀模式 `a,ba,b a`（去掉空格后为
`a,ba,ba`）；文本为空。多测时直接重新构造 \texttt{AC}，不要在旧对象上再次
\texttt{build()}。
\end{card}

\section{后缀自动机 SAM}

\cardmeta{高风险模板}{不同子串、子串出现次数、最长重复子串}{构建 \(O(n)\)，状态数 \(O(n)\)}{字符串长度决定节点池；每组数据必须 init}

\begin{lstlisting}[style=cpp]
struct SAM {
    struct Node {
        // go：从该状态出发的字符转移；len：该状态能表示的最长串长度
        // link：后缀链接；occ：在线插入时该状态被作为前缀经过的次数
        array<int, 26> go{};
        int link = 0, len = 0, occ = 0;
    };
    vector<Node> st;
    int last;

    void init(int n) {
        // SAM 最多 2n-1 个状态；1 号状态作为根，0 号位置不用
        st.assign(2 * n + 5, Node{});
        st[1] = Node{};
        last = 1;
        st.resize(2);
    }
    void extend(int c) {
        // 新建 cur，代表把字符 c 接到当前整串末尾
        int cur = st.size();
        st.push_back(Node{});
        st[cur].len = st[last].len + 1;
        st[cur].occ = 1;
        int p = last;
        while (p && !st[p].go[c]) {
            // 沿后缀链接补上所有缺失的 c 转移
            st[p].go[c] = cur;
            p = st[p].link;
        }
        if (!p) st[cur].link = 1;
        else {
            int q = st[p].go[c];
            if (st[p].len + 1 == st[q].len) st[cur].link = q;
            else {
                int clone = st.size();
                // q 的长度过长时复制出 clone，保持每个状态的长度区间合法
                st.push_back(st[q]);
                st[clone].len = st[p].len + 1;
                st[clone].occ = 0;
                while (p && st[p].go[c] == q) {
                    st[p].go[c] = clone;
                    p = st[p].link;
                }
                st[q].link = st[cur].link = clone;
            }
        }
        last = cur;
    }
    void count_occurrence() {
        // 按 len 降序把出现次数从子状态汇总到后缀链接父状态
        vector<int> ord(st.size() - 1);
        iota(ord.begin(), ord.end(), 1);
        sort(ord.begin(), ord.end(),
             [&](int a, int b) { return st[a].len > st[b].len; });
        for (int u : ord)
            if (st[u].link) st[st[u].link].occ += st[u].occ;
    }
};
\end{lstlisting}

\begin{card}
\textbf{完整题例：}P3804 的输入是一行字符串，要求所有子串中
\(\text{长度}\times\text{出现次数}\) 的最大值。把每个字符调用一次
\texttt{extend(s[i]-'a')}，再调用 \texttt{count\_occurrence()}，输出
\(\max_{u\ne 1}\texttt{len[u] * occ[u]}\)。

\begin{lstlisting}
输入：ababa
输出：6
\end{lstlisting}

\textbf{四组边界：}空串（答案 0）、`a`（1）、`aaaa`（6）、`ab`（2）。
若改问不同子串数量，答案是
\(\sum_{u\ne 1}(\texttt{len[u]}-\texttt{len[link[u]]})\)，不需要汇总
\texttt{occ}。每组数据必须重新调用 \texttt{init(n)}。
\end{card}

\section{后缀数组与 LCP}

\cardmeta{高风险模板}{后缀排序、LCP、不同子串}{排序版 \(O(n\log^2 n)\)}{工程上先用排序版，数据很大再换基数排序}

\begin{lstlisting}[style=cpp]
vector<int> suffix_array(string s) {
    int n = s.size();
    if (n == 0) return {};
    vector<int> sa(n), rk(n), tmp(n);
    iota(sa.begin(), sa.end(), 0);
    // rk[i] 是当前按前 2k 个字符比较时，后缀 i 的等价类编号
    for (int i = 0; i < n; ++i) rk[i] = s[i];
    for (int k = 1;; k <<= 1) {
        // 每轮按 (前半段排名, 后半段排名) 排序，倍增比较长度
        sort(sa.begin(), sa.end(), [&](int i, int j) {
            if (rk[i] != rk[j]) return rk[i] < rk[j];
            int ri = i + k < n ? rk[i + k] : -1;
            int rj = j + k < n ? rk[j + k] : -1;
            return ri < rj;
        });
        tmp[sa[0]] = 0;
        for (int i = 1; i < n; ++i) {
            int a = sa[i - 1], b = sa[i];
            int a2 = a + k < n ? rk[a + k] : -1;
            int b2 = b + k < n ? rk[b + k] : -1;
            tmp[b] = tmp[a] + (rk[a] != rk[b] || a2 != b2);
        }
        rk.swap(tmp);
        // 所有后缀排名都不同，排序完成
        if (rk[sa.back()] == n - 1) break;
    }
    return sa;
}

vector<int> kasai_lcp(const string &s, const vector<int> &sa) {
    int n = s.size();
    vector<int> rank(n), lcp(n);
    for (int i = 0; i < n; ++i) rank[sa[i]] = i;
    int h = 0;
    for (int i = 0; i < n; ++i) {
        int r = rank[i];
        if (r == 0) continue;
        int j = sa[r - 1];
        // h 是上一位置的 LCP，向右移动后最多只需减一。
        while (i + h < n && j + h < n && s[i + h] == s[j + h]) ++h;
        lcp[r] = h;
        if (h) --h;
    }
    return lcp; // lcp[r] = LCP(sa[r-1], sa[r])，lcp[0] = 0
}
\end{lstlisting}

\begin{card}
\textbf{完整题例：}字符串 `banana` 的后缀数组为
`5 3 1 0 4 2`，\texttt{kasai\_lcp} 返回 `0 1 3 0 0 2`。
\texttt{lcp[r]} 的下标跟后缀数组排名一致，任意两个后缀的 LCP 可继续在
\texttt{lcp} 上做 RMQ。

\textbf{四组边界：}空串；`a`；`aaaa`（大量相同前缀）；`ab`（无公共前缀）。
排序版复杂度为 \(O(n\log^2 n)\)，若 \(n\) 很大再替换排序器，不要在现场混用两版下标。
\end{card}

\section{回文自动机 PAM}

\cardmeta{高风险模板}{不同回文子串、最长回文后缀、出现次数}{构建 \(O(n)\)，节点数 \(\le n+2\)}{只支持小写字母；每组数据调用 \texttt{init}}

每个节点代表一个不同的回文串。0 号根长度为 \(-1\)，1 号根是空串；
\texttt{fail[u]} 指向去掉两端后最长的回文后缀。插入字符后，
\texttt{last} 是当前前缀的最长回文后缀。

\begin{lstlisting}[style=cpp]
struct PAM {
    struct Node {
        array<int, 26> go{};
        int fail = 0, len = 0;
        i64 occ = 0; // 先统计最长回文后缀次数，再沿 fail 汇总
    };
    vector<Node> tr;
    vector<int> text;
    int last = 1, n = 0;

    void init(int max_len) {
        tr.clear(); tr.reserve(max_len + 2);
        text.clear(); text.reserve(max_len);
        tr.push_back(Node{}); tr.push_back(Node{});
        tr[0].len = -1; tr[0].fail = 0; // -1 根
        tr[1].len = 0;  tr[1].fail = 0;  // 空串根
        last = 1; n = 0;
    }
    int get_fail(int u) const {
        while (n - 1 - tr[u].len < 0 ||
               text[n - 1 - tr[u].len] != text[n])
            u = tr[u].fail;
        return u;
    }
    int add(char ch) {
        int c = ch - 'a';
        text.push_back(c);
        int cur = get_fail(last);
        if (!tr[cur].go[c]) {
            int now = tr[cur].go[c] = tr.size();
            tr.push_back(Node{});
            tr[now].len = tr[cur].len + 2;
            tr[now].fail = (tr[now].len == 1)
                ? 1 : tr[get_fail(tr[cur].fail)].go[c];
        }
        last = tr[cur].go[c];
        ++tr[last].occ;
        ++n;
        return last;
    }
    void count_occurrence() {
        vector<int> ord(tr.size() - 2);
        iota(ord.begin(), ord.end(), 2);
        sort(ord.begin(), ord.end(),
             [&](int a, int b) { return tr[a].len > tr[b].len; });
        for (int u : ord) tr[tr[u].fail].occ += tr[u].occ;
    }
};
\end{lstlisting}

\begin{card}
\textbf{完整题例：}插入 `abacaba` 后，\texttt{tr.size()-2} 为 7，表示
`a,b,c,aba,bac,cab,abacaba` 等不同非空回文子串的数量。
输入 `a`、`aaaa`、`abacaba`、`ababa` 时输出分别为 `1,4,7,5`。
如果题目问出现次数，先调用 \texttt{count\_occurrence()}；如果只问不同回文串数量，
直接输出 \texttt{tr.size()-2}。多测必须重新 \texttt{init(max\_len)}。
\end{card}
